% LaTeX report template
% File: LaTeX_report_template.tex
% Usage notes:
% - Recommended engines:
%     * pdflatex + biber (or)
%     * xelatex + biber
% - If you use the 'minted' package for syntax highlighting, compile with -shell-escape.
%   Example: pdflatex -shell-escape main.tex
% - To compile bibliography with biblatex/biber: 
%     pdflatex main.tex
%     biber main
%     pdflatex main.tex
%     pdflatex main.tex

\documentclass[11pt,a4paper]{article}

% Language and encoding
\usepackage[utf8]{inputenc}
\usepackage[T1]{fontenc}
\usepackage[french]{babel}

% Page geometry and micro-typography
\usepackage[a4paper,margin=2.5cm]{geometry}
\usepackage{microtype}

% Math
\usepackage{amsmath,amssymb,amsthm}

% Graphics and floats
\usepackage{graphicx}
\usepackage{caption}
\usepackage{subcaption}
\usepackage{booktabs}
\usepackage{float}

% Code listings: prefer minted (requires -shell-escape), fallback to listings
\usepackage{ifthen}
\newboolean{useminted}
% Set to true if you want minted (remember -shell-escape). Set to false to use listings.
\setboolean{useminted}{false}
\ifthenelse{\boolean{useminted}}{%
  \usepackage{minted}
}{%
  \usepackage{listings}
  \lstset{basicstyle=\ttfamily\small,breaklines=true,frame=single}
}

% Bibliography with biblatex
\usepackage[backend=biber,style=numeric,sorting=nyt]{biblatex}
\addbibresource{refs.bib} % put your .bib file next to this tex file

% Hyperlinks
\usepackage[colorlinks=true,linkcolor=black,citecolor=black,urlcolor=blue]{hyperref}

% Header/Footer
\usepackage{fancyhdr}
\pagestyle{fancy}
\fancyhf{}
\lhead{\textsc{Nom du cours / Module}}
\rhead{\textsc{Rapport projet}}
\cfoot{\thepage}

% Title formatting
\usepackage{titling}
\pretitle{\begin{center}\LARGE\bfseries}
\posttitle{\par\end{center}\vskip 0.5em}
\preauthor{\begin{center}\large}
\postauthor{\end{center}}
\predate{\begin{center}\small}
\postdate{\end{center}}

% Theorem environment example
\newtheorem{theorem}{Théorème}[section]
\newtheorem{definition}[theorem]{Définition}

% Custom commands
\newcommand{\HRule}{\rule{\linewidth}{0.5mm}}

% DOCUMENT
\begin{document}

% Title page
\begin{titlepage}
  \centering
  \vspace*{2cm}
  {\scshape\LARGE Université / École \par}
  \vspace{1cm}
  {\huge\bfseries Titre du rapport \par}
  \vspace{1cm}
  {\Large Sous-titre (le cas échéant) \par}
  \vspace{2cm}
  \begin{tabular}{rl}
  \textbf{Auteur:} & Hugo \\
  \textbf{Encadrant:} & Prénom Nom \\
  \textbf{Année:} & 2025 \\
  \end{tabular}
  \vfill
  \vspace{0.8cm}
  {\small À remettre pour l'UE XXX}\
\end{titlepage}

% Abstract
\begin{abstract}
  Ici le résumé du rapport (100–250 mots). Expliquez brièvement le problème, la méthode et les résultats.
\end{abstract}

% Table of contents
\tableofcontents
\clearpage

% Sections
\section{Introduction}
Présentez le contexte, la problématique et l'organisation du rapport.

\section{Travaux connexes}
Faites une courte revue de la littérature, ou des sources principales.

\section{Méthodologie}
Expliquez la méthode, l'implémentation, les outils utilisés.

\subsection{Algorithmes}
Donnez des détails techniques, équations, pseudocode.

\begin{theorem}
Exemple de théorème important.
\end{theorem}

\section{Résultats}
Présentez vos résultats avec figures et tableaux.

\begin{figure}[H]
  \centering
  %\includegraphics[width=0.7\textwidth]{figure_exemple.pdf}
  \caption{Légende de la figure.}
  \label{fig:ex}
\end{figure}

\begin{table}[H]
  \centering
  \begin{tabular}{lrr}
    \toprule
    Méthode & Précision & Temps (s) \\
    \midrule
    Méthode A & 0.92 & 12.3 \\
    Méthode B & 0.88 & 8.5 \\
    \bottomrule
  \end{tabular}
  \caption{Exemple de tableau de résultats.}
  \label{tab:res}
\end{table}

\section{Discussion}
Analyse critique des résultats, limites, pistes d'amélioration.

\section{Conclusion}
Résumé et perspectives.

% Code example (minted or listings)
\section{Annexe technique — extrait de code}
\ifthenelse{\boolean{useminted}}{%
  \begin{minted}[fontsize=\small]{cpp}
  // Extrait C++ exemple
  #include <iostream>
  int main() {
    std::cout << "Bonjour monde" << std::endl;
    return 0;
  }
  \end{minted}
}{%
  \begin{lstlisting}[language=C++]
  // Extrait C++ exemple
  #include <iostream>
  int main() {
    std::cout << "Bonjour monde" << std::endl;
    return 0;
  }
  \end{lstlisting}
}

% Bibliography
\clearpage
\printbibliography[title={Références}]

% Appendices
\appendix
\section{Annexe A}
Matériel supplémentaire, données brutes, etc.

\end{document}
